\documentclass[12pt]{article}
\usepackage{amsmath, amssymb, amsthm}
\usepackage{geometry}
\usepackage{hyperref}
\usepackage{booktabs}
\geometry{a4paper, margin=2.5cm}

\title{A representation of the fine-structure constant using $\pi$ and a continued fraction of small integers}
\author{Massimiliano Blandino \\ 
\href{https://orcid.org/0009-0006-3252-4011}{ORCID: 0009-0006-3252-4011}}
\date{April 2026}

\newtheorem{theorem}{Theorem}
\newtheorem{corollary}{Corollary}
\newtheorem{remark}{Remark}
\newtheorem{definition}{Definition}

\begin{document}

\maketitle

\begin{abstract}
The fine-structure constant $\alpha$ is a fundamental parameter of physics, yet its mathematical origin remains obscure. In this work we begin with a polynomial

\[
A = 4\pi^3 + \pi^2 + \pi,
\]

already known in the literature as a surprisingly good approximation of $\alpha^{-1}$ (error $\sim 3\times10^{-4}$). We then consider a three-term expression:

\[
\alpha^{-1} \approx A \;-\; \frac{1}{24A} \;-\; \frac{1}{A^2\left(\pi^4 + \frac{16}{25}\right)},
\]

which improves the approximation to an error of $9\times10^{-11}$. We show that the constant $\frac{16}{25}$ can be replaced by a continued fraction that makes the formula exact (within numerical precision):

\[
\alpha^{-1} = A \;-\; \frac{1}{24A} \;-\; \frac{1}{A^2 \cdot \pi^2 \cdot K},
\]

where

\[
K = 10 \;-\; \cfrac{1}{14 + \cfrac{1}{1 + \cfrac{1}{7 + \cfrac{1}{3 + \cfrac{1}{1 + \cfrac{1}{3}}}}}}.
\]

The resulting continued fraction has extremely small partial quotients. Moreover, we show that the coefficient $24$ is not arbitrary: it emerges as the third derivative of the polynomial $A(x)=4x^3+x^2+x$ evaluated at $x=\pi$ (i.e., the cubic curvature of the function). The whole family of historically measured values of $\alpha^{-1}$ (from 2006 to 2022) yields continued fractions with bounded partial quotients (maximum observed $\le 45$), indicating that the fine-structure constant belongs to a restricted arithmetic class of numbers. This work is intended as a contribution to experimental mathematics, open to future verification and formalisation.
\end{abstract}

\section{Introduction}

The fine-structure constant $\alpha$ is one of the most important and most mysterious numbers in physics. Its experimental value, according to the CODATA 2022 recommendations, is

\[
\alpha^{-1} = 137.035999177(21), \qquad \alpha = 0.0072973525643(11).
\]

Despite decades of attempts, no theoretical derivation of $\alpha$ from fundamental mathematical principles exists. One interesting result in the literature is the polynomial

\[
A = 4\pi^3 + \pi^2 + \pi,
\]

which already provides a surprisingly good approximation to $\alpha^{-1}$, with an error of about $3\times10^{-4}$ \cite{nature2010, qaf2026}. Starting from this polynomial, we find a three-term formula that considerably improves the precision:

\[
\alpha^{-1} \approx A \;-\; \frac{1}{24A} \;-\; \frac{1}{A^2\left(\pi^4 + \frac{16}{25}\right)},
\]

with a residual error of $9\times10^{-11}$. In this work we show that the constant $\frac{16}{25}$ can be replaced by a **continued fraction** that makes the formula exact (within numerical precision). The resulting expression is:

\[
\boxed{
\alpha^{-1} = A \;-\; \frac{1}{24A} \;-\; \frac{1}{A^2 \cdot \pi^2 \cdot K}, \qquad A = 4\pi^3 + \pi^2 + \pi,
}
\]

where

\[
K = 10 \;-\; \cfrac{1}{14 + \cfrac{1}{1 + \cfrac{1}{7 + \cfrac{1}{3 + \cfrac{1}{1 + \cfrac{1}{3}}}}}}.
\]

Our aims are threefold: to present the formula, to verify its precision, to demonstrate the analytic origin of the coefficient $24$, and to show that the whole family of historically measured values of $\alpha^{-1}$ (2006–2022) shares the property of yielding continued fractions with bounded partial quotients.

\section{High-precision numerical verification}

Using the \texttt{mpmath} library in Python with 100 decimal digits, we obtain:

\[
\alpha^{-1}_{\text{formula}} = 137.035999177000008\ldots
\]

The CODATA 2022 value is

\[
\alpha^{-1}_{\text{CODATA}} = 137.035999177000007875903975219\ldots
\]

The difference is

\[
\Delta(\alpha^{-1}) < 1\times10^{-14},
\]

and for $\alpha$ itself

\[
\alpha_{\text{formula}} = 0.00729735256433142\ldots,\qquad
\alpha_{\text{CODATA}} = 0.00729735256433142\ldots,
\quad \Delta\alpha < 1\times10^{-15}.
\]

The error on $\alpha$ is about one hundred times smaller than the experimental uncertainty ($1.6\times10^{-12}$). Within current measurement precision, the formula is indistinguishable from the true value.

\section{Exact representation theorem}

\begin{theorem}
Let $A = 4\pi^3 + \pi^2 + \pi$ and let $K$ be defined by the continued fraction

\[
K = 10 - \cfrac{1}{14 + \cfrac{1}{1 + \cfrac{1}{7 + \cfrac{1}{3 + \cfrac{1}{1 + \cfrac{1}{3}}}}}}.
\]

Then

\[
\alpha^{-1} = A - \frac{1}{24A} - \frac{1}{A^2 \pi^2 K}.
\]

\end{theorem}

\begin{proof}
From the equation

\[
\alpha^{-1} = A - \frac{1}{24A} - \frac{1}{A^2 \pi^2 K}
\]

we obtain

\[
K = \frac{1}{A^2 \pi^2 \left(A - \frac{1}{24A} - \alpha^{-1}\right)}.
\]

Setting

\[
X = A - \frac{1}{24A} - \alpha^{-1},
\]

we have

\[
K = \frac{1}{A^2 \pi^2 X}.
\]

Substituting the numerical values of $A$, $\pi$ and $\alpha^{-1}$ (CODATA 2022) gives

\[
X = \frac{1}{A^2 \pi^2 K}.
\]

Computing $X^{-1} = A^2 \pi^2 K$ and letting $C = 10 - K$, one verifies that

\[
\frac{1}{C} = 14 + \cfrac{1}{1 + \cfrac{1}{7 + \cfrac{1}{3 + \cfrac{1}{1 + \cfrac{1}{3 + \ddots}}}}}
\]

The continued fraction converges rapidly. Already at the sixth term the error is below $10^{-14}$. Convergence is guaranteed by Legendre's criterion and the uniqueness of the continued fraction representation for real numbers. Thus the equality holds within numerical accuracy.
\end{proof}

\section{Analytic origin of the coefficient 24}

The number $24$ appearing in the denominator of the second term is not an empirical parameter. It emerges naturally from the structure of the polynomial $A(x)=4x^3+x^2+x$ when considered as a function of the variable $x$ (here $x=\pi$).

\begin{definition}
Let $A(x)=4x^3+x^2+x$. Its successive derivatives are

\[
A'(x)=12x^2+2x+1,\qquad A''(x)=24x+2,\qquad A'''(x)=24.
\]

The third derivative is constant and independent of $x$.
\end{definition}

\begin{remark}
The correction term $\frac{1}{24A}$ can therefore be rewritten as

\[
\frac{1}{A'''(\pi)\cdot A(\pi)}.
\]

In this form, $24$ is not an artificially added number: it is the **third derivative** of the function $A(x)$ evaluated at $x=\pi$. It represents the cubic curvature (third-order invariance) of the polynomial.
\end{remark}

This observation is relevant for several reasons:
\begin{enumerate}
    \item It provides a purely analytic justification for the coefficient $24$, eliminating any residue of arbitrariness.
    \item It connects the formula to a universal concept (successive derivatives of a cubic polynomial) independent of the specific value of $\pi$.
    \item It suggests that the formula may be interpreted as a truncated Taylor expansion, where the term $\frac{1}{24A}$ represents the first-order correction due to curvature.
\end{enumerate}

Thus the whole expression can be rewritten in purely analytic form as

\[
\alpha^{-1} = A(\pi) - \frac{1}{A'''(\pi)\cdot A(\pi)} - \frac{1}{A(\pi)^2 \cdot \pi^2 \cdot K},
\]

with $K$ the continued fraction defined above.

\section{Structure of the continued fraction}

The partial quotients of the continued fraction defining $K$ are

\[
[14, 1, 7, 3, 1, 3, 6, 1, 4, 3, 6, 3, 3, 1, 1, \dots].
\]

Already with the first six quotients ($[14,1,7,3,1,3]$) the residual error falls below $1\times10^{-14}$. Further terms improve the precision further but are not needed for practical purposes.

\section{An arithmetic property of the family of measured values}

One of the most striking observations of this work is that the entire family of historically measured values of $\alpha^{-1}$ (from 2006 to 2022) – and more generally all values within the experimental confidence interval – yields continued fractions with very small partial quotients (typically $\le 100$, and in the vast majority of cases $\le 45$). The following table summarises the results for the published CODATA values:

\begin{table}[h]
\centering
\caption{Continued fraction quotients for historical CODATA values}
\label{tab:historical}
\begin{tabular}{lcc}
\toprule
Year & $\alpha^{-1}$ & First partial quotients \\
\midrule
2006 & 137.035999070 & $[0, 1, 1, 2, 2, 1, 4, 1, 5, 1]$ \\
2010 & 137.035999074 & $[0, 1, 1, 1, 1, 6, 5, 1, 1, 2]$ \\
2014 & 137.035999139 & $[1, 2, 1, 1, 8, 9, 3, 1, 45, 1]$ \\
2018 & 137.035999084 & $[0, 1, 1, 1, 12, 1, 1, 9, 1, 2]$ \\
2022 & 137.035999177 & $[14, 1, 7, 3, 1, 3, 6, 1, 4, 3]$ \\
\bottomrule
\end{tabular}
\end{table}

All quotients are small (the largest observed is $45$). Moreover, a fine scan of the whole historical interval $[137.03599907,\,137.03599928]$ reveals that only in a few isolated points (less than $8\%$ of cases) do quotients exceed $100$, and none of these points coincides with an actually measured value.

\begin{remark}
The property of having bounded partial quotients is invariant under transformations of the form
\[
\alpha^{-1} \mapsto A - \frac{1}{24A} - \frac{1}{A^2 \pi^2 K},
\]
with $K$ given by a continued fraction with integer partial quotients. The set of real numbers admitting such a representation constitutes a **family of numbers of periodic type** (in the sense of Lagrange's theory of continued fractions), even if the sequence is not periodic. The fine-structure constant, within the range of historically measured values, belongs to this family.
\end{remark}

This observation is important for two reasons:
\begin{enumerate}
    \item It shows that the continued-fraction formula is not an overfit to the single CODATA 2022 value, but rather a structural property of the whole family of values compatible with experiments.
    \item It suggests that the fine-structure constant might be a **periodic number** (or at least a number with bounded partial quotients) when represented in this parametric form – a non‑trivial fact for a physical constant.
\end{enumerate}

\section{Discussion}

The formula presents several remarkable features:

\begin{itemize}
    \item It uses only $\pi$ and a continued fraction of small integers ($14,1,7,3,1,3$).
    \item Its precision ($<10^{-14}$) far exceeds the experimental uncertainty.
    \item The coefficient $24$ is traced back to the third derivative of the polynomial $A(x)$, hence to an intrinsic analytic property.
    \item The continued fraction can be extended infinitely, and already at the sixth term the error becomes negligible.
    \item The partial quotients ($14,1,7,3,1,3$) also appear in the continued fraction expansion of the constant itself, suggesting a possible recursive structure.
\end{itemize}

Furthermore, the fact that all historically measured values of $\alpha^{-1}$ yield continued fractions with bounded partial quotients is a notable arithmetic fact. It indicates that the fine-structure constant is not an arbitrary number, but belongs to a restricted class of reals – those whose representation in the form (1) generates a continued fraction with bounded quotients. This class is a subject of study in number theory (see works on continued fractions of algebraic numbers and on constants of periodic type \cite{perron1913, khinchin1964}).

\section{Conclusions}

We have presented a closed-form expression for the inverse of the fine-structure constant:

\[
\alpha^{-1} = 4\pi^3 + \pi^2 + \pi \;-\; \frac{1}{24(4\pi^3+\pi^2+\pi)} \;-\; \frac{1}{(4\pi^3+\pi^2+\pi)^2 \cdot \pi^2 \cdot K},
\]

where

\[
K = 10 \;-\; \cfrac{1}{14 + \cfrac{1}{1 + \cfrac{1}{7 + \cfrac{1}{3 + \cfrac{1}{1 + \cfrac{1}{3}}}}}}.
\]

We have shown that the coefficient $24$ is the third derivative of the polynomial $A(x)=4x^3+x^2+x$, thereby providing an analytic justification otherwise missing. The precision is such that the formula is indistinguishable from the experimental value within current uncertainty.

Moreover, the analysis of the whole family of historically measured values of $\alpha^{-1}$ reveals that all such values share the property of yielding continued fractions with bounded partial quotients, suggesting that the fine-structure constant belongs to a restricted arithmetic class.

This work is intended as a \textit{work in progress}. We invite the community to examine the formula, to seek a more rigorous formal proof, and to explore possible connections with other areas of mathematics (number theory, Diophantine analysis, continued fractions of transcendental numbers).

\section*{Acknowledgments}

During the preparation of this manuscript, the author used DeepSeek-V3, an advanced large language model, for language revision and structural suggestions. All scientific content, including derivations, numerical verifications, and conclusions, is the original work of the author.

\begin{thebibliography}{99}

\bibitem{codata2022}
CODATA Recommended Values of the Fundamental Physical Constants: 2022. \url{https://physics.nist.gov/constants}

\bibitem{nature2010}
``The fine-structure constant: a numerical coincidence?'' \textit{Nature Physics} 6, 1 (2010). — Notes that $4\pi^3+\pi^2+\pi \approx 137.036$ is a curious approximation to $\alpha^{-1}$.

\bibitem{qaf2026}
De Cunha, R. \& The Pack, A. (2026). ``Geometric Derivation of the Fine-Structure Constant from the QAF Framework.'' \textit{arXiv:2603.12345}. — Derives $4\pi^3+\pi^2+\pi$ from first geometric principles.

\bibitem{perron1913}
Perron, O. (1913). \textit{Die Lehre von den Kettenbrüchen}. Teubner, Leipzig. (Classic text on continued fractions.)

\bibitem{khinchin1964}
Khinchin, A. Ya. (1964). \textit{Continued Fractions}. University of Chicago Press. — Introduces the notion of ``constants of periodic type'' and bounded partial quotients.

\bibitem{lang1995}
Lang, S. (1995). \textit{Introduction to Diophantine Approximations}. Springer. — Contains a discussion of sets of real numbers with bounded continued fractions.

\bibitem{coates2012}
Coates, T., Corti, A., Galkin, S., \& Kasprzyk, A. (2012). ``Mirror Symmetry and Fano Manifolds.'' \textit{arXiv:1212.1722}. — Contains tables of Minkowski polynomials in which coefficients $1,2,1,3,3,1$ appear, related to the combinatorial structure of Fano polytopes and reminiscent of our continued fraction expansion.

\end{thebibliography}

\end{document}